\chapter{Introduction}

\section{Motivation}
\label{sec:motivation}

Regular expressions are a fundamental formalism for describing regular
languages and are among the most widely used abstractions in computer
science. They underpin applications including lexical analysis, input
validation, log processing, search tools, and network security. Across
these applications, the problem is regular expression matching. Given
a regular expression and an input string, determine whether the string
belongs to the language denoted by the expression. Formally, this is
the \emph{membership problem}.

For example, the expression \texttt{(ab|a)*} matches \texttt{"aba"},
as it decomposes into \texttt{"ab"} followed by \texttt{"a"}, but does
not match \texttt{"abb"}, since no sequence of subexpressions can
produce it.

For many applications, a Boolean answer is enough. An input validator
may only need to accept or reject a password or URL. 
Others require more. For instance, a lexer must determine how
input decomposes into tokens. Capturing groups, schema validation,
and information extraction require identifying which input parts
correspond to parts of the expression. Recognition alone is not
enough; the matcher must also recover the \emph{structure}.


Efficiently solving these problems remains an active research area.
The choice of algorithm influences performance, memory consumption, and
security. Once matching becomes parsing, the way in which match
structure is recovered also influences correctness and raises the
question of how to handle ambiguity when multiple matches are possible.

Existing approaches include backtracking and finite-automata-based
techniques, such as nondeterministic (NFA) and deterministic (DFA)
automata. These approaches have different trade-offs in implementation
complexity, memory consumption, and worst-case performance. Backtracking
engines are flexible and widely used, but may exhibit exponential
running time on constructed inputs, leading to regular expression
denial-of-service (ReDoS) vulnerabilities~\citep{cox2007}.
Automata-based approaches provide stronger worst-case guarantees but may
require preprocessing or suffer from state-space explosion.

Derivative-based algorithms provide an alternative by transforming the
regular expression as input symbols are consumed. Each derivative
represents the language remaining after an input symbol has been
consumed. By maintaining a correspondence between an expression
and its successive residuals, derivative-based methods provide an elegant
foundation for regular expression matching and, with suitable
extensions, structured parsing.

However, derivative-based techniques are less common in practical
regular expression tooling. Their use for practical matching, parsing,
and diagnostic tools motivates this thesis: "From Theory to Tools".




\section{Problem Statement}
\label{sec:problem-statement}

Derivative-based regular expression matching has been studied extensively
from a theoretical perspective. Brzozowski introduced regular expression
derivatives to construct a finite-state automaton directly from a
regular expression~\citep{brzozowski1964}, while Antimirov introduced partial
derivatives, which represent sets of residual expressions and are closely
related to nondeterministic automata~\citep{antimirov1996}. Building on
these ideas, Sulzmann and Lu developed a derivative-based algorithm for
POSIX-correct parsing and provided a Haskell reference
implementation~\citep{sulzmannlu2014}.

Sulzmann and Lu's work therefore provides a basis for structured
regular expression matching, but leaves several practical questions open.
In particular, these constructions can be investigated in a systems
programming language with respect to execution time and memory consumption.
Rust provides an appropriate setting through its ownership model and
fine-grained control over data representation and memory allocation.

Additionally, the partial-derivative-based greedy parser provides an
alternative for practical comparison with the derivative-based parser and
the bit-coded parse-tree representation enables investigation of its
effect on memory consumption.

The central problem addressed by this thesis is therefore to investigate
the practical behaviour of these derivative-based parsing constructions:
first, by implementing them in Rust and evaluating their performance and
resource usage; second, by comparing the resulting parsers with two
established regular expression engines. In addition, the thesis
investigates whether the intermediate derivative computations can be
exposed to provide useful explanations of successful and unsuccessful
matches.



\section{Contributions}
\label{sec:contributions}

This thesis makes the following contributions:

\begin{itemize}
  \item \textbf{Rust implementation of derivative-based parsers.}
  The thesis implements the parsing constructions developed by Sulzmann
  and Lu~\citep{sulzmannlu2014} in Rust, using their Haskell
  implementation as a reference. Both the derivative-based POSIX and
  partial-derivative-based greedy parsers are implemented in direct and
  bit-coded variants, resulting in four parser variants.

  \item \textbf{Empirical evaluation of the parser variants.}
  The four implementations are evaluated with respect to execution time
  and memory consumption using synthetic benchmarks and a corpus of
  real-world regular expressions from Snort/Suricata, SpamAssassin, and
  RegexLib, following the source categories investigated by Mamouras
  et al.~\citep{mamouras2024}. The derivative-based implementations are
  additionally compared with Rust's \texttt{regex} crate and Google's
  RE2.

  \item \textbf{Derivative-tracing diagnostics.}
  The implementation records and exposes intermediate derivative
  computations during matching. These traces show how the regular
  expression evolves as input symbols are consumed and form the basis
  of a diagnostic system for explaining successful and unsuccessful
  matches.
\end{itemize}




\section{Thesis Outline}
\label{sec:thesis-outline}

The remainder of this thesis follows the development of the parsing
constructions from their underlying formal concepts through their Rust
implementation to their empirical evaluation. Rather than separating
theoretical concepts and implementation into distinct chapters, the
relevant concepts are introduced alongside their implementation. 
This structure follows the development of the constructions by
Sulzmann and Lu while presenting their implementation and evaluation
in Rust.


\begin{itemize}
\item \textbf{Chapter~\ref{sec:regular-expressions} - Regular Expressions.}
This chapter introduces the syntax and semantics of regular expressions,
together with the Rust representation of their syntax used throughout
the implementation.

\item \textbf{Chapter~\ref{sec:derivatives} - Derivatives.}
This chapter introduces Brzozowski's derivatives and Antimirov's partial
derivatives as algorithms for deciding language membership, together
with their Rust implementation.

\item \textbf{Chapter~\ref{sec:matching-vs-parsing} - Matching vs. Parsing.}
Before applying these algorithms to parsing, this chapter distinguishes
matching from parsing and introduces the disambiguation problem that
arises when a parse tree, rather than only a Boolean result, is required.

\item \textbf{Chapter~\ref{sec:derivative-parser} - Derivative Parser.}
This chapter presents Sulzmann and Lu's derivative-based POSIX parser
~\citep{sulzmannlu2014}, which provides the main algorithmic basis for
this thesis, together with its Rust implementation in both direct and
bit-coded forms.

\item \textbf{Chapter~\ref{sec:partial-derivative-parser} - Partial-Derivative Parser.}
This chapter presents Sulzmann and Lu's greedy parser based on
Antimirov's partial derivatives, together with its direct and bit-coded
Rust implementations.

\item \textbf{Chapter~\ref{sec:frontend} - Front-End.}
This chapter describes how a regular expression string is parsed into
an external AST and translated into the internal representation used
by the parsers.

\item \textbf{Chapter~\ref{sec:evaluation} - Evaluation.}
With all four parser variants in place, this chapter evaluates their
correctness and performance with respect to execution time and memory
consumption. The evaluation uses both synthetic benchmarks and a corpus
of real-world regular expressions and compares the parser variants with
Rust's \texttt{regex} crate and Google's RE2.

\item \textbf{Chapter~\ref{sec:discussion} - Discussion.}
This chapter discusses and interprets the evaluation results, situates
the work relative to related approaches, and reflects on the translation
of the reference constructions from Haskell to Rust. It also discusses
the derivative-tracing system, including its usefulness and limitations.

\item \textbf{Chapter~\ref{sec:conclusion} - Conclusion.}
Finally, this chapter summarizes the main findings and contributions in
relation to the research questions introduced in this chapter. It also
discusses directions for future work, including potential extensions
and improvements to the parser implementations and diagnostic tooling.
\end{itemize}